\SourceLeaf{195}ren könnender Gebrauch zu machen ist. Vielleicht werde ich, wenn ich einmal
mehr Musse gewinne, als in meinen gegenwärtigen Verhältnissen, selbst
in Zukunft meine Untersuchungen bekannt machen.

Mit Hochachtung verharre ich
\begin{center}
\begin{tabular}{@{}l@{\qquad}r@{}}
 & Ewr. Wohlgeboren\\
Göttingen den 8. November & ergebenster Diener\\
\multicolumn{1}{c}{1824.} & \textsc{C. F. Gauss.}
\end{tabular}
\end{center}
\Needspace{9\baselineskip}
\medskip

\begin{center}
\textsc{Gauss} an \textsc{Olbers}. Göttingen, 3. Mai 1827.
\end{center}

.\ .\ .\ .\ .\ .\ .\ Vor etwa 6 Wochen hatte ich das Vergnügen, Hrn. \textsc{Dirichlet},
von dem ich, glaube ich, Ihnen schon einmal geschrieben habe, hier persönlich
kennen zu lernen. Ich erwähnte gegen ihn des schlechten Aufsatzes von
\textsc{Ivory}; er kannte ihn nicht selbst, sagte mir aber, Hr. \textsc{Fourier}
habe ihm gesagt, dass er »Unsinn« sei, und dass er (F.) sehr über einen
lobhudelnden Artikel darüber im Férussac gelacht habe. Ungefähr eben so wie
ich habe er auch über sein Mémoire von der Gleichgewichtsgestalt einer
rotirenden homogenen Flüssigkeit geurtheilt. Es war mir überraschend noch in
mehrern andern Beziehungen, über Personen und Sachen, eine ausserordentliche
Übereinstimmung Hrn. \textsc{Fouriers} mit meinem Urtheile zu erfahren; z.\,B.
über imaginäre Grössen, über die Unbeweisbarkeit der Geometrie a priori u. s.
w. Nachdem ich Hrn. \textsc{Dirichlet} z.\,B. über letztere meine Ansicht kurz
angedeutet hatte, sagte er, dass ihm \textsc{Fourier} die seinige fast mit den
nemlichen Worten gesagt habe. 
.\ .\ .\ .\ .\ .

